% =====================================================================
\section{The scalar face of $H1$: one identity, four exact models, one accounting
error, and twelve more families}
\label{ch:scalar-h1}
% =====================================================================

This chapter reports the wave that stopped writing $H1$ as a family of inequalities and asked
what single object the family names. The answer is available in the algebra already on the page:
the exact split \eqref{eq:split} of Chapter~\ref{sec:wgeo}, read as a definition rather than as a
decomposition, says that $H1(\delta)$ \emph{is} the scalar positivity $G''(0)\ge\delta$, with
$\delta_{\max}=G''(0)$. Four models whose conditioned curvature is exactly computable then decide
it: one of them fails, at $\beta>\beta^*=4.30661$, and the mechanism of failure is exact Savvidy
concavity. Comparing $\beta^*$ with the deconfinement couplings raised a question the campaign had
not asked --- whether $H1$ is even the right target for $\mathrm{SU}(3)$ --- and the wave withdrew
it: on the $\mathrm{SU}(3)$ lift of the same model $\delta$ \emph{increases} with $\beta$, so
$V_0(3)<0$, $\beta^*(3)$ does not exist, and $H1$ covers the whole $\mathrm{SU}(3)$ scaling region
\LE{}. What stays open is the step from that one-plaquette (patch) model to the theory \LM{}.
The same wave audited the five conditions (they are not five) and added families 46--57 to the
immune map, including the two structural statements that survive: $\kappa_W=0$ forbids monomials
and constrains no exponential, and the loop-equation principal symbol is never elliptic. One
further corner --- the polymer log-convexity route --- is closed quantitatively, $0.679<1.40$.
\textbf{No gap is proved and no new route is opened; Clay remains $0\%$.}

\subsection{One identity instead of a family}
\label{sec:scalar-identity}

Recall the exact split of Chapter~\ref{sec:wgeo}, \eqref{eq:split}:
$G''(0)=E[H_{vv}+\Phi''\mid\xi=0]-\operatorname{Var}(G_v+\Phi'\mid\xi=0)
=\Delta_{\mathrm{dyn}}+W_{\mathrm{geo}}+X$, with $\Delta_{\mathrm{dyn}}=E[H]-\operatorname{Var}(G)$
and $X=-2\operatorname{Cov}(G,\Phi')$. Subtracting $E[H]+W_{\mathrm{full}}$ from both sides,
$W_{\mathrm{full}}:=W_{\mathrm{geo}}+X$, gives an identity with no content beyond that display:
\begin{equation}
\underbrace{\operatorname{Var}_\mu\bigl(\langle v,\nabla S\rangle\mid0\bigr)}_{\operatorname{Var}}
-\underbrace{E_\mu[H_{vv}\mid0]}_{E[H]}-W_{\mathrm{full}}
\ =\ -\,G''(0).
\label{eq:scalar-identity}
\end{equation}
It uses no model, no uniformity, and no hypothesis that a gap exists. Substituting the curvature
form \eqref{eq:opencell-curv} of Chapter~\ref{ch:gap} --- which is
$\operatorname{Var}\le E[H]+W_{\mathrm{full}}-\delta$ --- and moving terms again yields the
decision rule
\begin{equation}
H1(\delta)\iff G''(0)\ge\delta
\qquad\text{(pointwise, mode by mode, for every }(\beta,U')\text{)},
\qquad
\boxed{\ \delta_{\max}=G''(0).\ }
\label{eq:scalar-rule}
\end{equation}

\begin{claimbox}{cT}
\claimhead{Theorem}{(the scalar face of $H1$)} \emph{Under the definitions of
Chapter~\ref{sec:wgeo}, \eqref{eq:scalar-identity} and \eqref{eq:scalar-rule} are algebraic
identities. Consequently:}
\begin{enumerate}
\item \emph{(total order)} \emph{the family $\{H1(\delta)\}_\delta$ is totally ordered by
  implication, $H1(\delta)\Rightarrow H1(\delta')$ for $\delta'\le\delta$; the family carries
  exactly the information of one number per coarse mode, $G''(0)$. ``$H1$'' without a $\delta$ is
  the statement $G''(0)>0$.}
\item \emph{(no independent variance inequality)} \emph{$\operatorname{Var}$, $E[H]$ and
  $W_{\mathrm{full}}$ are not separately constrained by $H1$ beyond their appearance in
  \eqref{eq:scalar-identity}; the sign lives in the curvature. A proposal that bounds
  $\operatorname{Var}$ without controlling the sign of $G''(0)$ is a claim about $G''(0)$ in
  disguise, not an independent route.}
\item \emph{(what a $\beta$-uniform route must deliver)} \emph{a route claiming $H1$ for all
  $\beta$ and all vacua must produce $G''(0)\ge\delta$ on the whole $(\beta,\text{vacuum})$
  plane; $\S\ref{sec:scalar-plaquette}$ exhibits a three-dimensional model where this fails
  exactly.}
\end{enumerate}
\LT{}
\end{claimbox}

\src{\code{ch05-theorems.tex} \eqref{eq:split} and \S\ref{sec:wgeo}; \code{ch08-immune.tex} \S1
(``the display is then $G''\ge\delta$ in curvature form''); \code{ch04-gap.tex}
\eqref{eq:opencell-curv}; \code{research/v3/species/H1A-EXACT.md} \S1.3 (identity, $\delta_{\max}$,
and the same reading of the corpus).}

\textbf{What this does not do.} It does not make $H1$ easier, and nothing here is a step toward
the gap. It relocates the difficulty into a single number: the open content is now (i) the sign and
size of $G''(0)$ at the physical point and (ii) its $(L,a)$-uniformity, and (ii) is the part that
is not closed \LM. Dimensional discipline is unchanged and is what keeps the two apart: the
curvature floor is a lattice-unit object, $G''_a=a^2m_{\mathrm{phys}}^2$, so every $\delta$ in this
chapter is a dimensionless curvature floor and none is a physical mass bound
(\code{ch03-existence.tex} \S\ref{sec:dim}, \eqref{eq:dim}) \LT.

\textbf{Corollary for route classes} \LN. Statement 3 above is not rhetoric: it is decided by the
next subsection, and it says that any proposal whose target is ``$H1$ for all $\beta$'' --- as
opposed to the paper's fixed-$\beta$, $(L,a)$-uniform, per-coarse-mode statement
(\code{ch13-machine.tex}:27--29) --- is dead by \emph{target}, independent of technique. What
survives is the fixed-coupling reading, and the honest way to state it is
\eqref{eq:scalar-rule}: prove $G''(0)\ge a^2\Delta_0^2$ uniformly, at the physical point, on a
gapped vacuum.

\subsection{Exact failure on a single plaquette}
\label{sec:scalar-plaquette}

The model is the one-plaquette $\mathrm{SU}(2)$ theory in the flat chart with a fixed boundary
holonomy: $U=e^{\xi}$, $\xi\in\mathbb R^3$, $|\xi|<2\pi$, $B=e^{i\theta\sigma_3/2}$, and
\begin{equation}
S=\beta\Bigl(1-\cos\tfrac{\rho}{2}\cos\tfrac{\theta}{2}
+\sin\tfrac{\rho}{2}\sin\tfrac{\theta}{2}\,u\Bigr),\qquad u=\xi_3/\rho,
\label{eq:scalar-plaq-action}
\end{equation}
with measure $e^{-S}J\,d^3\xi$, $J=\sin^2(\rho/2)/\rho^2$, $\Phi=-\log J$; take $v=e_3$ and
condition on $\langle v,\xi\rangle=a_0=0$. On that slice ($u=0$, $\rho'=0$) two exact
simplifications hold: $\Phi$ is radial, so $\Phi'\equiv0$ and therefore
$\operatorname{Var}(\Phi')=0$, $X=0$; and $\rho'=0$, $u''=0$ give
$H_{vv}=\beta c\sin(r/2)/(2r)$ with $c=\cos(\theta/2)$.

\begin{claimbox}{cT}
\claimhead{Theorem}{(exact one-plaquette curvature)} \emph{At $\theta=\pi$: $c=0$, so $E[H]\equiv0$
exactly; on the slice $S\equiv\beta$ independently of $r$, so the conditioned measure is exactly
Haar; hence}
\begin{equation}
\delta(\beta)=G''(0)=W_0-\beta^2V_0
=0.2051153-0.0110593\,\beta^2,
\qquad
\beta^*=\sqrt{W_0/V_0}=4.30661 ,
\label{eq:scalar-plaq-delta}
\end{equation}
\emph{with $V_0:=\operatorname{Var}_{\mathrm{Haar}}\bigl(\sin(r/2)/r\bigr)=0.01105929$. Therefore}
\begin{equation}
\delta(\beta)<0\iff\beta>\beta^*=4.30661 .
\label{eq:scalar-plaq-zero}
\end{equation}
\emph{The coefficient $W_0=1/(2\,\mathrm{Cin}(2\pi))=0.20511530$ is exact
(\eqref{eq:v3anchor} of Chapter~\ref{ch:uvface}); $V_0$ is a Haar variance computed by
deterministic quadrature and read digit for digit against the archive.}
\LT{}/\LE{}
\end{claimbox}

\src{\code{research/v3/species/H1A-EXACT.md} \S2.1 (model, simplifications, and the table below);
$W_0$ from \code{ch15-uvface.tex} \eqref{eq:v3anchor}; $V_0$ and the seven point values verified
against \code{research/v3/species/w7} \code{[P1]} and \code{[A0]} rows.}

\footnotesize
\begin{center}
\begin{tabular}{@{}rrrrr@{}}
\toprule
$\beta$ & $\operatorname{Var}$ & $E[H]$ & $W_{\mathrm{full}}$ &
$\delta_{\max}=\operatorname{Var}-E[H]-W_{\mathrm{full}}$\\
\midrule
$0$       & $0.000000$ & $0$ & $0.205115$ & $+0.205115$\\
$1$       & $0.011059$ & $0$ & $0.205115$ & $+0.194056$\\
$2.5$     & $0.069121$ & $0$ & $0.205115$ & $+0.135995$\\
$4.30661$ & $0.205115$ & $0$ & $0.205115$ & $0.000000$\\
$6$       & $0.398135$ & $0$ & $0.205115$ & $-0.193019$\\
$10$      & $1.105929$ & $0$ & $0.205115$ & $-0.900814$\\
\bottomrule
\end{tabular}
\end{center}

\normalsize
\textbf{External check} \LE. The failure is not a numerical artefact, and the numbers are not new:
five quantities agree with the corpus archive to the last printed digit
(Table~\ref{tab:scalar-archive}).

\footnotesize
\begin{center}
\begin{tabular}{@{}>{\raggedright\arraybackslash}p{6.0cm}rrl@{}}
\toprule
Quantity & This chapter & Archive & Relative difference\\
\midrule
$G''(\beta=0,\theta=\pi)$ & $0.205115$ & $0.20511530$ & $<10^{-6}$\\
$G''(\beta=2.5,\theta=\pi)$ & $0.135995$ & $0.135995$ & $<10^{-6}$\\
$G''(\beta=6,\theta=\pi)$ & $-0.193019$ & $-0.193019$ & $<10^{-6}$\\
$\operatorname{Var}(\beta=1)/\beta^2$ & $0.0110593$ & $0.01105929$ & $<10^{-5}$\\
$\beta^*$ & $4.30661$ & $4.30661$ & $0$\\
\bottomrule
\end{tabular}
\captionof{table}{Three-way digit-for-digit agreement between the recomputation of this chapter and
the archived one-plaquette values. The archive rows are the \code{[P1]} table and the \code{[A0]}
line of the one-plaquette run; agreement is the reason the failure below is called exact.}
\label{tab:scalar-archive}
\end{center}

\normalsize
\textbf{Why it fails: exact Savvidy concavity} \LT. At $\theta=\pi$ the action
\eqref{eq:scalar-plaq-action} is $S=\beta\bigl(1+\sin(\rho/2)u\bigr)$, linear in $u=\xi_3/\rho$;
the conditioned point $\xi_3=0$ is therefore \emph{not} a minimum --- the minimum sits at
$\xi_3=-\pi$. On the same slice $S$ does not depend on $r$, so the measure is exactly Haar and all
curvature comes from the Jacobian, while the $\beta^2$ variance term drives $G''$ down. Hence
``the curvature at the conditioned point'' \emph{is} ``the stability of that point'',
$\delta<0\iff$ the vacuum does not sit at the conditioned point. This is the exact scalar version of
the wall recorded in Chapter~\ref{sec:savvidy}: the Hessian $-D^2+2adF_B$ of the Savvidy cell is
indefinite (a negative transverse magnetic mode surviving the gauge quotient; the measured
negative fraction is $7.1\%$). What the solvable model adds is that the phenomenon is not an
artefact of the continuum or of large volume: three dimensions suffice, and the sign is exact.

\textbf{Scope of the refutation} \LN{}/\LM. Refuted: the reading ``$H1$ for all $\beta$, for all
vacua''. Not refuted: the paper's $H1$, which is fixed physical $\beta$ with $(L,a)$-uniformity and
a per-coarse-mode quantifier --- at the physical point $\beta_W=2.4$--$2.5$ one has
$\beta<\beta^*$ and $\delta=+0.136>0$. The consequence is a constraint on proof routes, not a
statement about Yang--Mills: any proof of $H1$ must carry vacuum-selection information (sign
control of $\Phi''$, hence of $W_{\mathrm{geo}}$); a proof that does not will be broken by this
three-dimensional model at $\beta>\beta^*$. Two honest limitations are recorded: the $\delta$ of
\eqref{eq:scalar-plaq-delta} is a lattice-unit curvature with no $a\to0$ scaling performed, and the
model is a single plaquette with a single boundary holonomy.

\textbf{A retracted computation, kept as a criterion.} The first version of this calculation gave
$G''(\beta=2.5)=+0.5427$: it double-counted a second-order term $s\cos(r/2)/(2r)$ of
$\sin(\rho/2)u$ which vanishes on the slice because $u''=0$ and $\rho'=0$. The criterion that
catches the error is not the final number but the archive's $E[H]\equiv0$, which forces any
candidate expression to vanish at $\theta=\pi$. Recorded so the error is not reused.

\subsection{Four solvable models: where the floor actually lives}
\label{sec:scalar-four}

Applying \eqref{eq:scalar-rule} to the models whose conditioned curvature is exact gives
Table~\ref{tab:scalar-models}. One model fails, one degenerates in the infrared, and two are
positive and uniform in the parameter they are supposed to be uniform in.

\footnotesize
\begin{longtable}{@{}>{\raggedright\arraybackslash}p{3.9cm}
  >{\raggedright\arraybackslash}p{3.1cm}
  >{\raggedright\arraybackslash}p{6.3cm}c@{}}
\caption{Four solvable models decided by \eqref{eq:scalar-rule}. $\delta$ is read off the identity,
so no separate variance inequality is proved or refuted in any row; the label records whether the
row's content is exact algebra (\LT) or computed (\LE).}\label{tab:scalar-models}\\
\toprule
Model & $\delta$ & Reading & Label\\
\midrule
\endfirsthead
\toprule
Model & $\delta$ & Reading & Label\\
\midrule
\endhead
\bottomrule
\endlastfoot
1. single plaquette $\mathrm{SU}(2)$, $\theta=\pi$, $\beta=2.5$ &
$+0.135995$ &
fails above $\beta^*=4.30661$; margin at the physical point is $0.136$ &
\LT/\LE\\
1$'$. same, $\beta=6$ &
$-0.193019$ &
$H1$ false; exact, not numerical, and the mechanism is Savvidy concavity &
\LT/\LE\\
1$''$. same, $\beta=\beta^*$ &
$0$ &
degenerate: exact zero, $\delta_{\max}=0$ &
\LT\\
1$'''$. same, $\theta=0$, any $\beta$ &
$\ge+0.2909$ &
true for every $\beta$, margin increasing in $\beta$; the vacuum sits at the conditioned point &
\LT/\LE\\
2. $2$D Yang--Mills (Migdal/heat kernel, exact) &
$\ge+2.0101$ &
true for all $\phi\in(0,\pi)$ and all $t=g^2A/2$; the minimum sits at $\phi=1.5$ and the value is
$t$-stable to $<0.3\%$ for $t\in[5,40]$ &
\LT/\LE\\
3. $P4$ periodic kernel, $\beta=2.5$ &
$+0.22196\pm0.00011$ &
$L$-uniform: positive intercept of the exponential fit, $q=0.849$; $\operatorname{Var}=E[H]=0$
exactly, so $\delta=W_{\mathrm{full}}$ &
\LT/\LE\\
4. $1$D Dirichlet string &
$\pi^2/L^2\to0$ &
degenerate: the ``mass'' is a finite-size effect $\propto1/L^2$, and at $L\to\infty$ $H1$ has no
content &
\LT\\
\end{longtable}

\normalsize
\textbf{The structural finding, which matters more than the seven numbers} \LT. In models 2, 3 and 4
the identities $\operatorname{Var}=0$ and $E[H]=0$ are exact, not small: the coarse mode is either
an exact flat direction of $S_W$ (the periodic chain, where every plaquette is a pure difference;
the $2$D class angle) or a Gaussian eigenmode whose force depends only on the conditioned
coordinate ($1$D Dirichlet, $G=\lambda\xi$). By \eqref{eq:scalar-identity} the whole content of
$H1$ in these models is therefore carried by the Haar/Jacobian term $W_{\mathrm{geo}}$. Two
consequences. First, this is the solvable version of the theorem \eqref{eq:wgeo}: the floor is
carried by Haar and Brascamp--Lieb, not by a rank-one interaction term. Second, and this is the
part to keep, \emph{a positive $\delta$ is a ceiling, not an atom}: at $\beta=0$ the periodic model
gives $G''=0.2688$ while $f_1^2=0$ (\code{ch15-uvface.tex}:236--241), which is precisely the
machine-checked statement \code{H1\_does\_not\_imply\_atom}. The variance half of \eqref{eq:scalar-identity}
is nontrivial in exactly one of the four models --- the non-flat, non-Gaussian one --- and that is
the one that fails.

\textbf{Two technical warnings the models supply} \LT. In $2$D Yang--Mills $\phi\to0$ is a
singularity of the marginal ($J(0)=0$, $G''\to+\infty$), so conditioning cannot be placed at the
$\mathrm{SU}(2)$ identity; the corpus's flat-chart anchor $1/(2\,\mathrm{Cin}(2\pi))$,
\eqref{eq:v3anchor}, is the correct base point and is exactly the $W_0$ of
\eqref{eq:scalar-plaq-delta}. And in $1$D the Dirichlet floor decays as $\pi^2/L^2$: the
Dirichlet-type class necessarily degenerates, which is why the corpus had to move to the periodic
kernel of model 3 --- whose intercept is positive with $q=0.849$, i.e.\ a genuine saturation
(\code{ch15-uvface.tex} \S\ref{ch:uvface-periodic}).

\subsection{$\beta^*$ against the critical coupling: the $\mathrm{SU}(3)$ alarm, and its
withdrawal}
\label{sec:scalar-beta-star}

The number $\beta^*=4.30661$ is exact inside a three-dimensional toy at fixed $\theta=\pi$, so
comparing it with a deconfinement coupling is not a theorem about Yang--Mills; it locates the
failure window of $H1$ relative to the scaling region, which is the question that decides whether
$H1$ is aimed at the right object. In the standard Wilson coupling $\beta_W=2N/g^2$ the
deconfinement points are
\begin{equation}
\beta_c\bigl(\mathrm{SU}(2)\bigr)=2.2986(6),\qquad
\beta_c\bigl(\mathrm{SU}(3)\bigr)=5.6924(2)
\label{eq:scalar-betac}
\end{equation}
(\code{research/v3/species/BETA-CONVENTION.md} \S2.2, which also fixes the convention: both are
$\beta_W$, not the SZZ $\beta_A$; sources are the reference estimate quoted by
Sale--Lucini--Giansiracusa, arXiv:2211.16273 \S4, and Bazavov--Berg, hep-lat/0701007 \S\,III,
$\beta_t(4)$; the in-paper occurrence is \code{ch13-machine.tex}:140) \LC.

\textbf{The $\mathrm{SU}(2)$ reading.} $\beta^*(2)=4.30661>\beta_c(\mathrm{SU}(2))=2.2986$: the
failure window $\beta>\beta^*$ lies \emph{above} the deconfinement coupling, so $H1$ covers the
entire scaling region $\beta_c\le\beta<\beta^*$, the physical point $2.5$ sits inside it at
$\delta=+0.136$, and the failure established in \S\ref{sec:scalar-plaquette} is a large-$\beta$
phenomenon outside $\beta_c$.

\textbf{The alarm, and why it was raised.} An earlier round of this writing
(\code{LOOP-RUNBOOK.md} \S11.39) raised the opposite worry, and it was a fair one: since
$\beta^*(2)=4.30661$ already lies \emph{below} $\beta_c(\mathrm{SU}(3))=5.6924$, a $\beta^*(N)$
growing slowly --- like $\sqrt N$, or not at all --- would put the whole $\mathrm{SU}(3)$ scaling
region inside the failure window, and $H1$ would then be a condition about the wrong region. That
objection is now \emph{withdrawn by measurement}, and the measurement runs the other way.

\begin{claimbox}{cE}
\claimhead{Evidence}{($\beta^*(3)$ does not exist)} \emph{On the canonical $\mathrm{SU}(N)$ lift of
the model of \S\ref{sec:scalar-plaquette} --- the Cartan-transverse slice, $S=\beta\bigl(1-
\operatorname{Re}\operatorname{Tr}(UB)/N\bigr)$, the de~Wit--'t~Hooft Jacobian, $B$ the centre
element, this being the lift for which $N=2$ reproduces
\eqref{eq:scalar-plaq-delta} --- a Monte Carlo run with $3\times10^{5}$ valid samples out of
$3.4\times10^{6}$ draws gives}
\begin{center}
\footnotesize
\begin{tabular}{@{}rrrr@{}}
\toprule
$\beta$ & $\delta$ & $E[H]$ & $\operatorname{Var}$\\
\midrule
$0$  & $+0.34368\pm0.0077$ & $0$ & $0$\\
$6$  & $+0.63393\pm0.0676$ & $0.3064$ & $0.03239$\\
$10$ & $+0.95757\pm0.1446$ & $0.6478$ & $0.07483$\\
\bottomrule
\end{tabular}
\end{center}
\emph{so $\delta$ is monotone \emph{increasing} in $\beta$; the $\beta^2$ coefficient of
$\delta(\beta)$ is therefore negative, $V_0(3)<0$ in the notation of
\eqref{eq:scalar-plaq-delta}, the opposite sign to the $\mathrm{SU}(2)$ value
$V_0=+0.0110593$. Hence $\beta^*(3)$ \textbf{does not exist} (no zero for $\beta\le10$), and since
$\beta_c(\mathrm{SU}(3))=5.6924\ll10$, $H1$ holds across the \emph{entire} $\mathrm{SU}(3)$
scaling region and beyond it. The $\mathrm{SU}(N)$ lift at $N=4$ points the same way ($\delta$ from
$+0.48678$ to $+0.99508$ through $\beta=18$), with errors too large to carry weight on their own.}
\LE{}
\end{claimbox}

\src{\code{research/unconditional/w7/betastar\_n34.log} (the $N=3$ and $N=4$ scans, the numbers
in the display verbatim); \code{research/v3/species/LOOP-RUNBOOK.md} \S11.46 (the reading
$V_0(3)<0$, the withdrawal of \S11.39, and the $\beta_c\ll10$ comparison); the closed form
$V_0=(N_2Z-N_1^2)/(4Z^2)$, $Z=\tfrac12\mathrm{Cin}\,2\pi$, which makes the $\mathrm{SU}(2)$
linear fit in $\beta^2$ exact rather than fitted (\code{LOOP-RUNBOOK.md} \S11.45).}

\textbf{What the withdrawal does and does not cover} \LE{}/\LM. It covers the objection exactly
where the objection was raised --- inside a solvable model at a fixed centre boundary holonomy ---
and it gives no part of $H1$. Three limits are recorded rather than hidden: the run is Monte Carlo
($3\times10^{5}$ valid samples, acceptance $0.0882$) and its errors grow with $\beta$
($\pm0.1446$ at $\beta=10$); it is a single plaquette, or patch, not the whole theory; and one
boundary holonomy has been tested. What is still missing is the step from the patch to the theory,
i.e.\ the $(L,a)$-uniform statement at the physical point, which is $H1$ itself and is the one
substantive target left in \S\ref{sec:scalar-five}.

\begin{claimbox}{cT}
\claimhead{Theorem}{(the target is aimed correctly, conditionally on the model)} \emph{In the
solvable model, and conditional on the measurement above, the answer to the subsection's question
is no: for neither group is $H1$ pointed at the wrong region. For $\mathrm{SU}(2)$ the one-plaquette
failure is confined to $\beta>\beta^*(2)=4.30661>\beta_c=2.2986$, outside the scaling region; for
$\mathrm{SU}(3)$ there is no $\beta^*$ at all, since $\delta$ grows with $\beta$, and
$\beta_c=5.6924$ lies well inside the region where $\delta>0$. The ``wrong target'' objection of
\S11.39 is therefore withdrawn.}
\LC{}
\end{claimbox}

\textbf{Why this does not touch the paper's $H1$} \LT. The paper's condition is fixed-$\beta$,
$(L,a)$-uniform, per coarse mode (\code{ch13-machine.tex}:27--29). The toy's failure at
$\beta>\beta^*$ cannot refute it --- at $2.5$ the toy is on the safe side --- and the $\mathrm{SU}(3)$
measurement cannot prove it either. What both do is narrow the space in which the objection could
live: the ``$\beta$-uniform'' phrasing of $H1$ is now known to fail only above $\mathrm{SU}(2)$'s
scaling region, not inside any scaling region that matters.

\subsection{The five conditions are not five}
\label{sec:scalar-five}

Chapter~\ref{sec:missing} stated the conditions the missing tool must satisfy. Read again, they do
not have the same logical type, and the corpus's practice of scoring them as five independent axes
is an accounting error.

\begin{claimbox}{cN}
\claimhead{No-go}{(the five-item scorecard)} \emph{Among the five conditions:}
\begin{enumerate}
\item \emph{$C5\Rightarrow C2$, so $C2$ is not an independent axis. $C5$ is exactly $C2$ plus the
  axis qualifications (uniform in $(L,a)$ and surviving the continuum limit), so every tool
  satisfying $C5$ satisfies $C2$ and not conversely. The paper states the collapse itself:
  ``These are not six problems; they are one --- the volume-uniform lower bound on
  $m_{0^{++}}$ along the path'' (\code{ch13-machine.tex}:49--50). The window coordinate of $C2$
  ($\rho^2\le C^2$) is a restatement of that same floor, not an independent half
  (\code{ch13-machine.tex}:494, 504).}
\item \emph{$C3$ and $C4$ are filters, not propositions.} \emph{Their text is ``must not assume''
  and ``must fail''; such a condition can never be proved, and never be refuted --- only violated
  by a candidate. Non-circularity is already folded into family 47 (presupposing the conclusion is
  a way to die), and failing the free field is admission question 2 of \S11.1. Scoring a filter
  pays for nothing.}
\item \emph{$C1$ arrives with $C5$ inside the setting} \LC: \emph{at $\beta\approx2.5$, $g\approx1.5$
  and there is no small parameter, so a bound uniform in $(L,a)$ there cannot be the product of a
  convergent expansion.}
\end{enumerate}
\emph{Consequently the scorecard that counts five axes --- ``$1.5/5$'', ``$2/5$'' and their
relatives --- manufactures free progress: every formal construction passes $C1$ (it does not
expand) and $C3$ (it presupposes nothing) automatically, so a score of $2/5$ can mean zero
content.}
\LN{}
\end{claimbox}

\src{\code{research/v3/species/H1B-FIVE.md} \S4.1--4.4 (the implication, the restatement, the
filters, and the list of scorecard occurrences); scorecard occurrences in
\code{BC-BETHE.md}:193, \code{LF-LASTFOUR.md}:79/130/184/273, \code{MG-META.md}:240/254/266,
\code{SDP-CERT.md}:295; the collapse sentence and the window remark in \code{ch13-machine.tex}
:49--50, :494, :504.}

\textbf{What is left: one target} \LM. After the collapse exactly one substantive object remains,
and it is the one the paper already marks: $C5$'s \emph{rate uniformisation}. What survives the
strong-coupling chain is
$\mathrm{RP}\Rightarrow\mathrm{KP}(a)\Rightarrow$ exponential clustering $\Rightarrow\mathrm{KL}$,
hence a gap $\ge a$ with $a$ in lattice units and no scale separation near the root
($\xi/a\approx0.5$--$1.3$); the tool that would turn that into an $(L,a)$-uniform statement is the
corpus object $d_5$/$(\mathrm{CP})$, marked \emph{untouched} in the paper
(\code{ch13-machine.tex}:171--173). Two remarks keep the target from being quietly widened. First,
replacing the single triangle inequality $\rho=N/2-C^*N\lvert\beta\rvert$ \emph{is} $H1$, not a
bypass of it (\code{LOOP-RUNBOOK.md}:368, 377). Second, the paper's own subsection ``The five
conditions, as they now stand'' covers only $C2$, $C3$ and $C4$ --- $C1$ and $C5$ are absent
(\code{ch13-machine.tex}:491--517) --- so the same omission appears twice, independently, and both
times at $C5$. The next wave has one problem, not five.

\subsection{Families 46--57 of the immune map}
\label{sec:scalar-families}

The immune map of Chapter~\ref{ch:immune} gained twelve families in this wave. A family is a
mechanism class: if the class dies, everything built out of it inherits the death, which is why the
list is worth keeping even where the individual proposals were uninteresting.
Table~\ref{tab:scalar-fam} gives one line each, with the number that carries the killing.

\footnotesize
\begin{longtable}{@{}>{\raggedright\arraybackslash}p{3.35cm}
  >{\raggedright\arraybackslash}p{11.05cm}@{}}
\caption{Immune-map families 46--57, one line each. Every row is a class-level no-go \LN{} except
row 55, which is an admission rule rather than a no-go; row 47 is the family that absorbs the
filter $C3$ of \S\ref{sec:scalar-five}.}\label{tab:scalar-fam}\\
\toprule
Family & What kills it, with the number that carries it\\
\midrule
\endfirsthead
\toprule
Family & What kills it, with the number that carries it\\
\midrule
\endhead
\bottomrule
\endlastfoot
46. BD/Polchinski--LSI import &
four directions, each fatal once: $L\to\infty$ fails on $L^*=16.66$; $a\to0$ fails because
$\gamma/\lambda_1=ma/2\to0$; $\beta\to\infty$ fails because $Q\propto\beta^{1.14}$; and the acid
test fails because $\gamma_{\mathrm{free}}=m^2$ makes the criterion return $N/4$\\
47. every shape or scaling argument &
any statement about the \emph{scaling} of $m$ is logically posterior to $m>0$: $m=C\Lambda$,
$C=6.3(3)$, slope $=$ one-loop coefficient all presuppose the conclusion --- same order gives a
tautology, mixed order gives a falsehood\\
48. numerical continuum extrapolation &
reaching $90\%$ of the asymptotic slope needs $\beta\approx27$ ($\xi/a\sim10^{11}$); at accessible
$\beta$ the slope is only $0.53$ ($\mathrm{SU}(3)$) and $0.21$ ($\mathrm{SU}(2)$), and at
$N_\tau=4$ its sign is reversed\\
49. averaging over $\nu$ (the $D3$ class) &
$c$ is lower semicontinuous on the Gibbs class and $\nu$ has full support, so
$\mathrm{ess}$-$\mathrm{sup}=\sup$ and $D3\equiv D2$: averaging weakens nothing\\
50. $H2$-global type extra hypotheses &
incomparable with $H1$; the conjunction gives $m=\lambda\lambda_f/(\lambda+\lambda_f)\le\min$, which
only loses --- provable at $\beta<1/48$ where it is not needed, and needed at $2.5$ where it is not
provable\\
51. $u$-completeness &
the ceiling is $\beta_c^{\mathrm{KP}}\le1.182$ and $r\equiv0$ for $\beta>1.0408$; the smallest
\emph{closed} polymer has $\lvert X\rvert=6$, so every estimate at $\lvert X\rvert=1$ describes the
empty set\\
52. SDP localisation hierarchies &
a hierarchy is legitimate iff its truncation is PSD-monotone; the measure and kernel have no such
structure (relative values $-3.3$ to $-8.0\times10^{-2}$), and the required frustration-free
common-kernel decomposition does not exist for $4$D Yang--Mills\\
53. TRG/EKR analytic-ball import &
the ball has radius $0.1587$ and the physical tensor leaves it at $0.65$, $0.876$, $1.17$ in three
norms --- $2.9$--$3.8\times$ --- and the radius is a property of the map, not of the computation\\
54. more compute, higher precision &
truncation stopped being the bottleneck long ago ($j\le1$ gives $4\times10^{-5}$ at $\beta=2.5$; the
$6j$ table to $J=7$ takes $1.03$ s; $d\le6\beta^2$ is a theorem), so the obstruction was never
precision\\
55. ``a new species'' as the answer &
the five paradigm classes (positivity, convexity, spectral, Perron, RG) are closed by the 44
families of Chapter~\ref{ch:immune}; a new species must name the class it escapes and why it is not
a re-dressing of it\\
56. the loop-algebra sector &
$\kappa_W=0$ forbids power laws; the area law solves the loop equation only in the degenerate
sector; every non-degenerate sector has a hyperbolic symbol (family 57); and the strongest
same-level conclusion is a floor, not a gap --- hence no gap\\
57. loop-space coordinates &
the principal symbol is never elliptic, $\det=-\tfrac{(\rho_1R_1-\rho_2R_2)^2}{16\pi^2R_1^2R_2^2}
\le0$, so every maximum-principle, Harnack, comparison and monotonicity tool is structurally
unavailable in this coordinate\\
\end{longtable}

\normalsize
Two of these deserve a paragraph, because they are the ones most likely to be misused.

\textbf{The weakening of $\kappa_W=0$} \LT{}/\LM. A corollary of loop-space scale covariance,
$\kappa_W=0$, is routinely quoted as if it constrained the loop profile. It does not.
With $\kappa_W(s):=d\ln W(sC)/d\ln s=\epsilon(sx)$ and $\epsilon(x)=xF'(x)/F(x)$, the quantity is
$s$-independent only for a pure monomial, where $\kappa_W=-p$; every stretched exponential gives
$\kappa_W(s)=-\alpha c(sx)^\alpha$, which \emph{depends on $s$}, so for that class the equation
``$\kappa_W=0$'' is not a constraint but a type error
(\code{KW1-CONS.md} \S1.2, with a four-profile table and a numerical kernel). Hence $\kappa_W=0$
forbids monomial prefactors --- zero anomalous dimension --- and says nothing about $\alpha$; read
pointwise it would exclude the area law $e^{-\sigma R^2}$ itself, which the corpus has as a
solution, and that contradiction is the reductio that shows the reading is wrong
(\code{KW1-CONS.md} \S2.2). The exponent $\alpha=2$ has an owner and it is not $\kappa_W$: it is
the $R$-power $q=1$ of the contact source in the additive sector, $\alpha=q+1$. So $\kappa_W=0$ is
a weak structural theorem and cannot carry the gap, which fixes the final wording of family 56:
$\kappa_W=0$ forbids power laws, the area law solves the equation only in the degenerate sector,
and the non-degenerate sectors have a hyperbolic symbol --- no gap
(\code{LOOP-RUNBOOK.md}:1665).

\textbf{The principal symbol is never elliptic} \LT. In the two-variable (``figure-eight'')
sector the principal symbol of the loop equation has
\begin{equation}
\det=a_{11}a_{22}-\tfrac14a_{12}^2
=-\frac{(\rho_1R_1-\rho_2R_2)^2}{16\pi^2R_1^2R_2^2}\ \le\ 0,
\label{eq:scalar-det}
\end{equation}
with equality only on a hypersurface, verified at five configurations including an exact parabolic
one (maximum deviation $3.5\times10^{-18}$) (\code{FIG8.md} \S3, eq.~(0.5.5);
\code{LOOP-RUNBOOK.md}:1546--1551). The symbol is hyperbolic, not elliptic, and every classical
tool of the maximum-principle family --- maximum principle, Harnack, comparison, monotonicity ---
requires ellipticity. The eight-line theorem of the species has every step of that type
($f''$ sign $\Rightarrow$ $f'$ monotone $\Rightarrow$ comparison of $f'$ with an integral
$\Rightarrow$ bounded $1/f$), so in this sector the mechanism is not \emph{inherited}: it cannot be
expressed. What this buys is a no-go, and it should be quoted as one: non-ellipticity does not
force decay either, so family 57 removes a tool class without supplying a replacement.

\subsection{A quantitative no-go for the polymer route}
\label{sec:scalar-polymer}

One corner of the arithmetic route was left open by the earlier rounds: prove that the per-area
decay rate of the surface weight beats the surface entropy,
$\lim_A\bigl(-\log\lvert w_\beta(A)\rvert\bigr)/A>\kappa$. Three statements now close it.

\textbf{(1) It is a restatement} \LN. With $w_\beta(A)\sim e^{-c(\beta)A}$ one has
$c(\beta)=\kappa+\sigma(\beta)a^2(\beta)/\pi$, hence $c>\kappa\iff\sigma>0$: proving $c>\kappa$ is
proving the area law. The corpus's own theorem proves both at once,
$\sigma=-\log(c_f/c_0)-\gamma$ with $\gamma$ the surface-counting entropy, i.e.\ with $\kappa$
itself (\code{research/confinement/PROOF.md}, theorem B, as quoted in
\code{Q4-LOGCONVEX.md} \S1.1, \S4).

\textbf{(2) The direction is wrong for the purpose} \LT. \emph{Direction lemma.} With $w_A>0$ and
$r_A:=-\log w_A/A$: if $A\mapsto\log w_A$ is convex then $r_A$ is monotone non-increasing, so
$\lim r_A=\inf_A r_A\le r_{A_0}$ and every finite-$A$ reading is an \emph{upper} bound; if
$A\mapsto\log w_A$ is concave then $r_A$ is monotone non-decreasing and finite-$A$ readings are
lower bounds. The criterion needs $\liminf r_A>\kappa$, i.e.\ a lower bound, so log-convexity can
only refute decay and can never prove it (\code{Q4-LOGCONVEX.md} \S3). The corpus's proposal also
attached premise and conclusion to opposite objects: ``$-\log w$ convex'' \emph{is} ``$w$
log-concave''.

\textbf{(3) At the physical point it is quantitatively false} \LE. Two corpus conventions give
$u(2.5)=1.014390$ and $0.507195$, hence $c(2.5)=-0.014288$ and $+0.678860$; the stronger of the
two, $0.679$, is still below $\kappa_{\min}=1.40$. And $\kappa$ is not small. With the correct
four-dimensional face degree $\Delta=4(2d-3)=20$ --- not $12$, which is the three-dimensional value;
brute-force enumeration gives $4,12,20,28$ for $d=2,\dots,5$ --- the encoding bound yields
$\mu\le4(\Delta-1)=76$, hence
\begin{equation}
\kappa\in[\ln4.06,\ \ln76]=[1.4012,\ 4.3307],
\label{eq:scalar-kappa}
\end{equation}
the lower end by $Z^2$ injection being \LC{} and the upper end \LT{}
(\code{Q5-KAPPA.md} \S2, \S4). So $0.679<1.4012\le\kappa$: entropy beats weight. The failure is
already visible at the smallest closed polymer --- the $\lvert X\rvert=6$ cube has weight
$6.809\times10^{-2}>0$ while $e^{\kappa A}=e^{8.4}=4.4\times10^{3}$, so $\Omega_Aw\approx3\times
10^{2}>1$ --- and $c(\beta)$ decreases with $\beta$ (it tends to $0$ as $\beta\to\infty$), so
$c>\kappa$ holds only for $\beta\lesssim1$ ($c=1.40$ at $\beta=0.4983$ in the first convention,
$1.0290$ in the second), which is the strong-coupling region, not $2.5$
(\code{Q4-LOGCONVEX.md} \S5).

\begin{claimbox}{cN}
\claimhead{No-go}{(polymer log-convexity)} \emph{The surviving corner of the arithmetic route is
dead three ways: it is a restatement of $\sigma>0$; its natural convexity runs in the direction
that yields upper bounds, not the lower bound the criterion needs; and at $\beta=2.5$ it is false
by a factor, $0.679<1.40$. What survives is the pair of bounds
\eqref{eq:scalar-kappa}, whose width $2.93$ cannot be narrowed by finite-$A$ enumeration: a width
below $1$ would need $\mu_{\max}<e\,\mu_{\min}\approx11.0$, and the enumerated
$a_5^{1/5}=11.43$ already crosses that. The route cannot be closed by sharpening $\kappa$.}
\LN{}
\end{claimbox}

\src{\code{research/v3/species/Q4-LOGCONVEX.md} \S1.1, \S3, \S4, \S5 (the equivalence, the
direction lemma, and the two conventions at $\beta=2.5$); \code{research/v3/species/Q5-KAPPA.md}
\S1--\S4 ($\Delta=20$ with the enumeration check, $\mu\le76$, and the width argument);
\code{research/v3/species/KW1-CONS.md} \S5 for the ``cannot carry the gap'' reading that keeps
this no-go from being read as a statement about confinement.}

\subsection{Verdict of this chapter}
\label{sec:scalar-verdict}

\begin{itemize}
\item \LT{} $H1$ has exactly one numeric content per coarse mode, $\delta_{\max}=G''(0)$
  \eqref{eq:scalar-rule}. The family of inequalities is an illusion of notation; what must be
  proved is a single scalar positivity, and it must be proved uniformly.
\item \LT{}/\LE{} The scalar positivity \emph{fails exactly} on one plaquette,
  $\delta(\beta)=0.2051153-0.0110593\beta^2<0\iff\beta>\beta^*=4.30661$, digit-for-digit against
  the archive, and the mechanism is exact Savvidy concavity: the curvature at the conditioned point
  is the stability of that point.
\item \LE{}/\LM{} The target is aimed correctly, in the model:
  $\beta^*(2)=4.30661>\beta_c(\mathrm{SU}(2))=2.2986$, so the one-plaquette failure sits above
  $\mathrm{SU}(2)$'s scaling region; and at $N=3$ the lift has $\delta$ \emph{increasing} in $\beta$
  ($+0.34368$, $+0.63393$, $+0.95757$ at $\beta=0,6,10$), so $V_0(3)<0$ and $\beta^*(3)$ does not
  exist --- $H1$ then holds across the entire $\mathrm{SU}(3)$ scaling region, $\beta_c=5.6924\ll10$.
  The ``wrong target'' alarm is withdrawn; what stays open is the lift from the patch model to the
  theory \LM.
\item \LN{}/\LM{} The five conditions collapse to one: $C5\Rightarrow C2$, $C3$ and $C4$ are
  filters, $C1$ comes with the setting. The only substantive target left is $C5$'s rate
  uniformisation (\code{ch13-machine.tex}:171--173), still untouched.
\item \LN{} Families 46--57 join the immune map, of which two are structural rather than
  proposal-specific: $\kappa_W=0$ forbids monomials and constrains no exponential, and
  \eqref{eq:scalar-det} says the loop-equation symbol is never elliptic, so the whole
  maximum-principle tool class is unavailable there.
\item \LN{}/\LE{} The polymer log-convexity route is quantitatively closed: $c(2.5)=0.679<\kappa$,
  and $\kappa\in[1.4012,4.3307]$ cannot be narrowed by the method that needs narrowing.
\end{itemize}

\textbf{Not done, and not to be read as done.} No $a\to0$ scaling is performed on any $\delta$ of
this chapter; the $\mathrm{SU}(2)$ model $2$ split $(\operatorname{Var},E[H])$ is fixed by the
identity rather than by independent quadrature; the constants of model $3$ are quoted from the
corpus rather than recomputed here; the step from the one-plaquette (patch) model to the theory is
missing at every $N$, including the $N=3$ sign that lifts the alarm; and the true value of
$\kappa$ is unknown between the bounds. None of these changes a sign in the tables above, and none
of them is a gap. \textbf{Clay is $0\%$: this chapter narrows the target and closes twelve families
plus one corner; it opens no route and proves no mass gap.}
